\documentclass[11pt,a4paper]{article}
\usepackage{turkos}

\begin{document}
\phaseheader{1}{x86 Assembly and the Boot Process}{Learn enough assembly to read and write kernel glue code, and follow a PC from power-on to your first instruction}{5}{10}

\ataglance{2--3 weeks}{3}{Phase 0 complete: cross toolchain, NASM, QEMU, GDB}{A NASM ``phrase book'' in your journal, two
assembly functions called from C, a 512-byte boot sector that prints a message, and a Multiboot kernel that both
\code{qemu -kernel} and a GRUB ISO boot. You will prove it runs by finding \code{0xCAFEBABE} in \code{EAX}.}

\section{Why this phase}

You don't need to become an assembly programmer to write an operating system. Turk-OS will contain a few hundred
lines of assembly in total. But those lines do the things C cannot do: set up the very first stack, load descriptor
tables, save and restore registers on an interrupt, switch between processes, and drop into user mode. If you can't
read them, you can't debug them. And the only way to see what is really happening when the kernel crashes is to
read the machine state in terms of registers, the stack and instructions.

The second half of the phase answers a question that feels like magic until you see it: how does a PC go from
power-on to running \emph{your} code? You will follow that path once by hand, writing a boot sector the BIOS loads
directly, and then hand the job to GRUB, which is what Turk-OS will use from here on.

\section{Concepts you need}

\subsection{Registers}

The 32-bit x86 CPU has eight general-purpose registers. Each can be used as a whole (32 bits) or, for the first
four, in smaller pieces that keep their names from the 8086 era (Figure~\ref{fig:eax}). The remaining registers
have fixed jobs, and several of them will get a phase of their own.

\begin{figure}[htbp]
\centering
\begin{tikzpicture}[x=0.36cm, y=1cm]
  \draw[draw=tkNavy!70, fill=tkGray!10] (0,0) rectangle (16,0.7);
  \draw[draw=tkNavy!70, fill=tkBlue!12] (16,0) rectangle (24,0.7);
  \draw[draw=tkNavy!70, fill=tkRed!10] (24,0) rectangle (32,0.7);
  \node[font=\sffamily\scriptsize, text=tkGray] at (8,0.35) {upper 16 bits (no separate name)};
  \node[font=\sffamily\small\bfseries] at (20,0.35) {AH};
  \node[font=\sffamily\small\bfseries] at (28,0.35) {AL};
  \foreach \x/\b in {0.5/31, 15.5/16, 16.5/15, 23.5/8, 24.5/7, 31.5/0}
    \node[font=\sffamily\scriptsize, text=tkGray] at (\x,0.95) {\b};
  \draw[decorate, decoration={brace, amplitude=5pt, mirror}, tkNavy, line width=0.6pt] (16,-0.1) -- (32,-0.1)
    node[midway, below=5pt, font=\sffamily\small\bfseries] {AX (16 bits)};
  \draw[decorate, decoration={brace, amplitude=6pt, mirror}, tkRed, line width=0.6pt] (0,-0.75) -- (32,-0.75)
    node[midway, below=6pt, font=\sffamily\small\bfseries, text=tkRed] {EAX (32 bits)};
\end{tikzpicture}
\caption{One register, four names. The same split applies to EBX, ECX and EDX.}
\label{fig:eax}
\end{figure}

\begin{center}
\small\renewcommand{\arraystretch}{1.25}
\begin{tabularx}{\linewidth}{@{}l >{\raggedright\arraybackslash}X >{\raggedright\arraybackslash}X@{}}
\toprule
\tkth{Register(s)} & \tkth{Job} & \tkth{Where Turk-OS uses it} \\
\midrule
EAX, EBX, ECX, EDX & general purpose; EAX holds function return values & everywhere; EAX also carries the system call number (Phase 10) \\
ESI, EDI & general purpose; source and destination for string instructions & copying memory \\
ESP & stack pointer: address of the top of the stack & set up before any C runs (Phase 2); swapped on every context switch (Phase 8) \\
EBP & frame pointer: base of the current function's stack frame & walking the stack for backtraces (Phase 3) \\
EIP & instruction pointer: address of the next instruction & printed in every crash report \\
EFLAGS & condition flags (zero, carry, sign...) and the interrupt flag IF & \code{cli}/\code{sti} (Phase 4) \\
CS, DS, SS, ES, FS, GS & segment selectors & loaded from your own GDT (Phase 3) \\
CR0, CR2, CR3, CR4 & control registers: protected mode, paging, fault address & paging (Phase 6) \\
\bottomrule
\end{tabularx}
\end{center}

\subsection{A small instruction vocabulary}

About twenty instructions cover nearly all the assembly in this project. NASM uses Intel syntax, where the
destination comes first: \code{mov eax, 5} means \code{eax = 5}. Square brackets mean ``the memory at this address'':
\code{mov eax, [ebx+4]} reads four bytes starting at address \code{ebx+4}. When the size is ambiguous, say it:
\code{mov dword [eax], 0}.

\begin{center}
\small\renewcommand{\arraystretch}{1.2}
\begin{tabularx}{\linewidth}{@{}l X l@{}}
\toprule
\tkth{Instruction} & \tkth{Meaning} & \tkth{C equivalent} \\
\midrule
\code{mov dst, src} & copy a value & \code{dst = src;} \\
\code{add}, \code{sub}, \code{inc}, \code{dec} & arithmetic; results set EFLAGS & \code{a += b; a--;} \\
\code{and}, \code{or}, \code{xor}, \code{not}, \code{shl}, \code{shr} & bitwise operations and shifts & \code{a &= m; a <<= 2;} \\
\code{cmp a, b} then \code{je}/\code{jne}/\code{jl}/\code{jg}/\code{jge} & compare (a subtraction that only sets flags), then jump on the result & \code{if (a == b) goto L;} \\
\code{test a, a} then \code{jz} & bitwise AND that only sets flags; common zero check & \code{if (a == 0) goto L;} \\
\code{jmp label} & unconditional jump & \code{goto label;} \\
\code{push x}, \code{pop x} & \code{esp -= 4; [esp] = x} and the reverse & (the call stack) \\
\code{call f}, \code{ret} & push the return address and jump; pop it and jump back & \code{f();} \code{return;} \\
\code{lea eax, [ebx+ecx*4]} & compute an address without reading memory & \code{eax = &ebx[ecx];} \\
\code{in al, dx}, \code{out dx, al} & read or write an I/O port & (device access, Phase 2) \\
\code{cli}, \code{sti}, \code{hlt} & disable or enable interrupts; halt until the next interrupt & (Phase 4) \\
\code{int n}, \code{iret} & software interrupt; return from an interrupt handler & (Phases 3 and 10) \\
\bottomrule
\end{tabularx}
\end{center}

\subsection{The stack and the cdecl calling convention}

C code compiled for 32-bit x86 follows the \textbf{cdecl} convention, and your assembly must follow it too whenever
it calls C or is called from C. The caller pushes the arguments \emph{right to left}, then executes \code{call},
which pushes the return address. Inside the callee, the first argument is therefore at \code{[esp+4]}, the second at
\code{[esp+8]}, and so on. The return value goes in EAX. The callee may freely overwrite EAX, ECX and EDX
(\emph{caller-saved}), but must restore EBX, ESI, EDI, EBP and ESP before returning (\emph{callee-saved}). The caller
removes the arguments afterwards. Most functions start by saving EBP and pointing it at their frame, so arguments
sit at fixed offsets from EBP (Figure~\ref{fig:frame}). Remember the callee-saved list: it is the reason the context
switch in Phase 8 only saves four registers.

\begin{figure}[htbp]
\centering
\begin{tikzpicture}[slot/.style={draw=tkNavy!70, minimum width=4.2cm, minimum height=0.62cm, font=\sffamily\small, fill=tkLight},
  lab/.style={font=\ttfamily\small, anchor=west}]
  \node[slot] (b) at (0,0) {argument \code{b}};
  \node[slot, below=0pt of b] (a) {argument \code{a}};
  \node[slot, below=0pt of a, fill=tkAmber!12] (ra) {return address};
  \node[slot, below=0pt of ra, fill=tkBlue!10] (bp) {saved EBP};
  \node[slot, below=0pt of bp] (l1) {local variable};
  \node[slot, below=0pt of l1, fill=white] (l2) {\dots};
  \node[lab] at ([xshift=4pt]b.east) {[ebp+12]};
  \node[lab] at ([xshift=4pt]a.east) {[ebp+8]};
  \node[lab] at ([xshift=4pt]ra.east) {[ebp+4]};
  \node[lab] at ([xshift=4pt]bp.east) {[ebp]  $\leftarrow$ EBP};
  \node[lab] at ([xshift=4pt]l1.east) {[ebp-4]};
  \node[lab] at ([xshift=4pt]l2.east) {$\leftarrow$ ESP};
  \draw[tkarrow] ([xshift=-10pt]b.north west) -- ([xshift=-10pt]l2.south west) node[midway, left, tknote, text width=1.8cm] {stack grows toward lower addresses};
  \draw[decorate, decoration={brace, amplitude=4pt}, tkGray] ([xshift=3.3cm]b.north east) -- ([xshift=3.3cm]a.south east) node[midway, right=5pt, tknote] {pushed by caller};
\end{tikzpicture}
\caption{Stack frame inside \code{f(a, b)} after \code{push ebp; mov ebp, esp}. Higher addresses are at the top.}
\label{fig:frame}
\end{figure}

\subsection{Real mode, protected mode, and how a PC boots}

An x86 CPU wakes up in \textbf{real mode}, pretending to be a 1978 8086: 16-bit registers, addresses formed as
\code{segment * 16 + offset}, one megabyte of reachable memory, no protection at all. The BIOS runs in real mode and
so does the boot sector it loads. Operating systems switch to \textbf{protected mode} as soon as they can: 32-bit
registers, a 4\,GiB address space, privilege levels, and (later) paging. GRUB does this switch for you. Figure
\ref{fig:boot} shows the whole chain.

\begin{figure}[htbp]
\centering
\begin{tikzpicture}[node distance=0.45cm and 0.45cm, every node/.style={font=\sffamily\scriptsize}]
  \node[tkbox, text width=2.3cm] (pw) {\textbf{Power on}\\CPU in real mode, starts at \code{0xFFFFFFF0}};
  \node[tkbox, text width=2.3cm, right=of pw] (bios) {\textbf{BIOS firmware}\\self-test, finds a bootable device};
  \node[tkbox, text width=2.5cm, right=of bios] (mbr) {\textbf{Boot sector}\\512 bytes loaded at \code{0x7C00}, must end in \code{0x55 0xAA}};
  \node[tkbox, text width=2.5cm, right=of mbr] (grub) {\textbf{GRUB}\\enters protected mode, reads the Multiboot header};
  \node[tkhi, text width=3.4cm, below=0.8cm of grub] (k) {\textbf{Your kernel} (\code{loader})\\loaded at 1\,MiB; \code{EAX = 0x2BADB002}, \code{EBX} = Multiboot info};
  \node[tkok, text width=3.2cm, left=1.2cm of k] (q) {\textbf{\code{qemu -kernel}}\\QEMU's built-in Multiboot loader: skips BIOS and GRUB};
  \draw[tkarrow] (pw) -- (bios); \draw[tkarrow] (bios) -- (mbr); \draw[tkarrow] (mbr) -- (grub);
  \draw[tkarrow] (grub) -- (k); \draw[tkarrow, dashed] (q) -- (k);
\end{tikzpicture}
\caption{From power-on to the kernel. In Task~1.3 you replace the boot sector with your own; from Task~1.4 on GRUB (or QEMU) does it.}
\label{fig:boot}
\end{figure}

The boot sector is loaded low in memory because the first megabyte is crowded (Figure~\ref{fig:lowmem}). Your kernel
goes at 1\,MiB, the first address above all of it.

\begin{figure}[htbp]
\centering
\begin{tikzpicture}[m/.style={draw=tkNavy!70, minimum width=6.2cm, minimum height=0.52cm, font=\sffamily\scriptsize, fill=tkLight},
  a/.style={font=\ttfamily\scriptsize, anchor=east}]
  \node[m, fill=tkGreen!10] (k) at (0,0) {extended memory: \textbf{Turk-OS kernel loaded here}};
  \node[m, below=0pt of k, fill=tkGray!15] (bios) {system BIOS ROM};
  \node[m, below=0pt of bios, fill=tkGray!15] (rom) {video BIOS and option ROMs};
  \node[m, below=0pt of rom, fill=tkAmber!12] (vga) {video memory (text buffer at \code{0xB8000})};
  \node[m, below=0pt of vga, fill=tkGray!15] (ebda) {extended BIOS data area (size varies)};
  \node[m, below=0pt of ebda] (free2) {free conventional memory};
  \node[m, below=0pt of free2, fill=tkRed!10] (bs) {\textbf{boot sector} (512 bytes)};
  \node[m, below=0pt of bs] (free1) {free conventional memory};
  \node[m, below=0pt of free1, fill=tkGray!15] (bda) {BIOS data area};
  \node[m, below=0pt of bda, fill=tkGray!15] (ivt) {real-mode interrupt vector table};
  \foreach \n/\t in {k/0x00100000, bios/0x000F0000, rom/0x000C0000, vga/0x000A0000, ebda/\textasciitilde 0x0009FC00, free2/0x00007E00, bs/0x00007C00, free1/0x00000500, bda/0x00000400, ivt/0x00000000}
    \node[a] at ([xshift=-4pt]\n.south west) {\t};
\end{tikzpicture}
\caption{The first megabyte of a PC's physical memory (not to scale). Addresses mark the bottom of each region.}
\label{fig:lowmem}
\end{figure}

\begin{conceptbox}{The Multiboot contract}
When GRUB or QEMU jumps to your kernel, the Multiboot specification (\S3.2, \emph{Machine state}) guarantees:
EAX holds the magic value \code{0x2BADB002}; EBX holds the physical address of the \emph{Multiboot information}
structure (memory size, memory map, boot modules); the CPU is in 32-bit protected mode with paging off; the segment
registers describe flat 4\,GiB segments; interrupts are disabled. It also warns you of two things it does \emph{not}
give you. \textbf{ESP is undefined}, so you must set up your own stack before calling any C function (Phase 2).
And \textbf{the GDTR may be invalid}, so you must not reload a segment register until you have installed your own
GDT (Phase 3). To be recognised, the kernel must contain a 12-byte Multiboot header (magic \code{0x1BADB002},
flags, checksum) aligned to 4 bytes within the first 8\,KiB of the file.
\end{conceptbox}

\section{Reading plan}

\begin{readingplan}
\must & \LOB & Ch.~2 \emph{First Steps}: Booting, Hello Cafebabe (compiling, linking, GRUB, ISO, running) & 12--17 & The exact steps of Tasks 1.4--1.5. \\
\must & \OSC & \S2.9 \emph{Building and Booting an Operating System} & 92--95 & Firmware, boot loaders and the boot sequence in general. \\
\should & \MOS & \S1.3.1 \emph{Processors}; \S1.3.6 \emph{Booting the Computer} & 21--24, 34--35 & Registers, pipelines and privileged mode; the BIOS boot path. \\
\should & \OSDI & \S2.6.6 \emph{Bootstrapping MINIX 3} & 156--160 & How a real x86 OS gets from the boot sector to its kernel. \\
\could & \OSDI & \S4.6.2 \emph{Segmentation with Paging: The Intel Pentium} & 415--420 & First look at real versus protected mode addressing. \\
\could & \OSC & \S9.6 \emph{Example: Intel 32- and 64-bit Architectures} & 379--383 & Preview of the address translation you will build in Phase 6. \\
\end{readingplan}

Three external references matter here. The \textbf{NASM manual}, chapters 2--3, covers syntax, labels and
directives. The \textbf{Multiboot Specification 0.6.96}, section 3, defines the header and the machine state above.
The \textbf{Intel SDM Volume 1}, chapter 3 (basic execution environment) and chapter 6 (procedure calls and the
stack), is the authoritative version of Figure~\ref{fig:frame}.

\begin{minixbox}
MINIX 3 does not use GRUB. It ships its own boot monitor, and \OSDI\,\S2.6.6 walks through what it does: load the
kernel and the system processes into memory, switch to protected mode, and jump to the kernel's assembly entry point
in \texttt{kernel/mpx386.s (06300)}. Skim the first page of that file now. You won't understand all of it yet, but
you will recognise the stack setup and the jump into C (\texttt{kernel/start.c (06900)}), which is exactly what you
build in Phase 2.
\end{minixbox}

\section{Build it}

\task{Learn to read compiler output}
Continue the phrase book you started in Phase 0. For each C construct below, write the smallest function that uses
it, compile with \code{i686-elf-gcc -O0 -S -masm=intel}, and annotate every instruction in your journal: a
local variable, an \code{if}, a \code{for} loop, an array index, a struct field access, a pointer dereference, and a
call with two arguments. Then compile the same code with \code{-O2} and notice what disappears.
\verify{For the two-argument call, you can find the two \code{push} instructions, the \code{call}, and the
\code{add esp, 8} that removes the arguments.}

\task{Write assembly that C can call}
Install the 32-bit C runtime on your host (\code{sudo dnf install glibc-devel.i686 libgcc.i686}) so you can build
32-bit Linux test programs. Then write two functions in NASM and call them from a C test with \code{assert}. The
first is a leaf function that reads its arguments straight from ESP. The second builds a proper frame and uses a
callee-saved register, so it must preserve it.

\begin{lst}{asm}{playground/funcs.s}
; int max(int a, int b)      cdecl: a at [esp+4], b at [esp+8], result in eax
global max
global sum_array
section .text
max:
    mov  eax, [esp + 4]     ; eax = a
    mov  ecx, [esp + 8]     ; ecx = b   (ecx is caller-saved: free to use)
    cmp  eax, ecx
    jge  .done              ; if (a >= b) return a
    mov  eax, ecx           ; else return b
.done:
    ret

; int sum_array(const int *a, int n)
sum_array:
    push ebp
    mov  ebp, esp           ; now [ebp+8] = a, [ebp+12] = n
    push ebx                ; ebx is callee-saved: we must restore it
    mov  ebx, [ebp + 8]     ; ebx = a
    mov  ecx, [ebp + 12]    ; ecx = n
    xor  eax, eax           ; sum = 0
.loop:
    test ecx, ecx
    jz   .done
    add  eax, [ebx]         ; sum += *a
    add  ebx, 4             ; a++   (an int is 4 bytes)
    dec  ecx
    jmp  .loop
.done:
    pop  ebx
    pop  ebp
    ret
\end{lst}

\begin{lst}{sh}{terminal}
nasm -f elf32 playground/funcs.s -o build/funcs.o
gcc -m32 -g playground/test_funcs.c build/funcs.o -o build/test_funcs && ./build/test_funcs
\end{lst}
\verify{The asserts pass. Then step through \code{sum_array} in GDB with \code{layout asm}, \code{stepi} and
\code{info registers}, and watch ESP change on every \code{push} and \code{pop}.}

\task{A boot sector that says hello}
This is a detour, and a valuable one. You will write the first 512 bytes a PC executes from disk, in 16-bit real
mode, using a BIOS service to print. Turk-OS itself will not use this code; GRUB replaces it in the next task.

\begin{lst}{asm}{playground/boot.s}
bits 16                     ; the CPU is in 16-bit real mode
org  0x7C00                 ; the BIOS loads us at this address

start:
    xor  ax, ax
    mov  ds, ax             ; data segment = 0, so [si] means address si
    mov  si, msg
.next:
    lodsb                   ; al = [ds:si], then si = si + 1
    test al, al
    jz   .halt              ; stop at the terminating zero
    mov  ah, 0x0E           ; BIOS "teletype output" function
    int  0x10               ; ask the BIOS video service to print al
    jmp  .next
.halt:
    cli
    hlt
    jmp  .halt

msg: db "Merhaba, Turk-OS!", 0

times 510 - ($ - $$) db 0   ; pad with zeros up to byte 510
dw 0xAA55                   ; boot signature: bytes 0x55, 0xAA
\end{lst}

\begin{lst}{sh}{terminal}
nasm -f bin playground/boot.s -o build/boot.bin
ls -l build/boot.bin             # exactly 512 bytes
hexdump -C build/boot.bin | tail -n 2   # last bytes: 55 aa
qemu-system-i386 -drive format=raw,file=build/boot.bin
\end{lst}
\verify{QEMU prints your message after the SeaBIOS banner. Delete the \code{dw 0xAA55} line and confirm that
the BIOS now refuses to boot the disk.}

\task{Hello Cafebabe: a Multiboot kernel}
Now the real Turk-OS entry point, following \LOB chapter 2. The kernel only loads a recognisable value into EAX and
spins, so you can prove your code is running by inspecting the register. Put the Multiboot header in its own section
and have the linker place it first, which guarantees it lands in the first 8\,KiB of the file.

\begin{lst}{asm}{boot/loader.s}
global loader                   ; the entry symbol, named in link.ld

MAGIC_NUMBER equ 0x1BADB002     ; tells GRUB this is a Multiboot kernel
FLAGS        equ 0x0            ; no extra requests yet (Phase 2 changes this)
CHECKSUM     equ -(MAGIC_NUMBER + FLAGS)   ; magic + flags + checksum == 0

section .multiboot
align 4
    dd MAGIC_NUMBER
    dd FLAGS
    dd CHECKSUM

section .text
loader:
    mov eax, 0xCAFEBABE         ; a value you will recognise in a register dump
.loop:
    jmp .loop                   ; spin forever
\end{lst}

\begin{lst}{plain}{link.ld}
ENTRY(loader)                   /* where execution starts */
SECTIONS
{
    . = 0x00100000;             /* load everything from 1 MiB upwards */

    .text ALIGN(0x1000) : {
        *(.multiboot)           /* Multiboot header first */
        *(.text)
    }
    .rodata ALIGN(0x1000) : { *(.rodata*) }
    .data   ALIGN(0x1000) : { *(.data) }
    .bss    ALIGN(0x1000) : { *(COMMON) *(.bss) }
}
\end{lst}

Build it with \code{nasm -f elf32 boot/loader.s -o build/loader.o}, then
\code{i686-elf-ld -T link.ld -o build/kernel.elf build/loader.o}. Run it with QEMU's built-in Multiboot loader and
open the monitor on your terminal:

\begin{lst}{sh}{terminal}
grub2-file --is-x86-multiboot build/kernel.elf && echo "multiboot header OK"
qemu-system-i386 -kernel build/kernel.elf -monitor stdio
(qemu) info registers
\end{lst}
\verify{\code{info registers} shows \code{EAX=cafebabe} and an \code{EIP} just above \code{0x00100000}.}

\task{Boot through a real GRUB ISO}
\LOB builds its ISO with GRUB Legacy's \code{stage2_eltorito}. On a current Fedora system, GRUB 2 does the whole job
with \code{grub2-mkrescue}. Create a GRUB configuration and let the tool assemble a bootable CD image around your
kernel:

\begin{lst}{plain}{boot/grub.cfg}
set timeout=0
set default=0

menuentry "Turk-OS" {
    multiboot /boot/kernel.elf
    boot
}
\end{lst}

\begin{lst}{sh}{terminal}
mkdir -p build/iso/boot/grub
cp build/kernel.elf build/iso/boot/
cp boot/grub.cfg build/iso/boot/grub/
grub2-mkrescue -o build/turkos.iso build/iso
qemu-system-i386 -cdrom build/turkos.iso -monitor stdio
\end{lst}
\verify{Same result as Task~1.4: \code{EAX=cafebabe}. You now boot exactly like a real PC would.}

\task{A Makefile for the kernel}
Typing these commands gets old fast. Write a Makefile with the targets \code{all}, \code{run}, \code{iso},
\code{debug} and \code{clean}. Remember that recipe lines must start with a real Tab character.

\begin{lst}{make}{Makefile}
AS    := nasm
LD    := i686-elf-ld
BUILD := build
OBJS  := $(BUILD)/loader.o

all: $(BUILD)/kernel.elf

$(BUILD)/kernel.elf: $(OBJS) link.ld
	$(LD) -T link.ld -o $@ $(OBJS)

$(BUILD)/%.o: boot/%.s
	@mkdir -p $(BUILD)
	$(AS) -f elf32 $< -o $@

run: $(BUILD)/kernel.elf
	qemu-system-i386 -kernel $< -monitor stdio

iso: $(BUILD)/kernel.elf boot/grub.cfg
	mkdir -p $(BUILD)/iso/boot/grub
	cp $< $(BUILD)/iso/boot/kernel.elf
	cp boot/grub.cfg $(BUILD)/iso/boot/grub/
	grub2-mkrescue -o $(BUILD)/turkos.iso $(BUILD)/iso

debug: $(BUILD)/kernel.elf
	qemu-system-i386 -kernel $< -s -S &
	gdb $< -ex "target remote localhost:1234" -ex "break loader" -ex "continue"

clean:
	rm -rf $(BUILD)

.PHONY: all run iso debug clean
\end{lst}
\verify{\code{make clean && make run} rebuilds and boots from scratch.}

\task{Your first kernel debugging session}
Run \code{make debug}. QEMU starts frozen (\code{-S}) and waits for GDB on port 1234 (\code{-s}). GDB stops at
\code{loader}. Try \code{info registers}, \code{stepi}, \code{x/4xw $esp} and \code{p/x $eax}, and use
\code{layout asm} for a live disassembly view. You will use this workflow in every phase from now on, so practise it
while the kernel is still two instructions long.
\verify{After one \code{stepi}, \code{p/x $eax} prints \code{0xcafebabe}.}

\section{Testing and debugging}

\begin{pitfalls}
QEMU or BIOS: ``No bootable device'' with your boot sector & File is not exactly 512 bytes, or the last two bytes are not \code{55 AA} & Check with \code{ls -l} and \code{hexdump -C}; keep the \code{times} padding line. \\
Boot sector prints garbage & Missing \code{org 0x7C00} or \code{bits 16}, or DS not zero & All label addresses depend on \code{org}; set DS before \code{lodsb}. \\
GRUB: ``no multiboot header found'' & Header not in the first 8\,KiB, not 4-byte aligned, or a bad checksum & Keep \code{*(.multiboot)} first in \code{.text}; run \code{grub2-file --is-x86-multiboot}. \\
QEMU \code{-kernel}: ``Error loading uncompressed kernel without PVH ELF Note'' & QEMU found no Multiboot header, or the ELF is 64-bit & Same checks; make sure you used \code{nasm -f elf32} and the \code{i686-elf} linker. \\
\code{objdump} output looks backwards & AT\&T syntax (source first) & Use \code{objdump -d -M intel} to match NASM. \\
Your C test crashes after calling \code{sum_array} & A callee-saved register (EBX) was not restored, or the stack is unbalanced & Every \code{push} needs a matching \code{pop} before \code{ret}. \\
\end{pitfalls}

\section{Milestone}

\begin{milestonebox}[The kernel boots]
\begin{checklist}
  \item \code{max} and \code{sum_array} in NASM pass their C tests, and you stepped through them in GDB.
  \item Your boot sector prints ``Merhaba, Turk-OS!'' in QEMU and is exactly 512 bytes.
  \item \code{kernel.elf} boots with \code{make run} and with the GRUB ISO; \code{EAX=0xCAFEBABE} in both.
  \item \code{make debug} stops at \code{loader} and you can single-step.
  \item You can draw the cdecl stack frame for a two-argument call from memory.
\end{checklist}
\end{milestonebox}

\begin{questionbox}
\begin{enumerate}
  \item Why is the boot sector loaded at \code{0x7C00}, and why must it end with \code{0x55 0xAA}?
  \item What does the Multiboot checksum guarantee? Compute it for \code{FLAGS = 0x3}.
  \item Right after \code{call f}, what is at \code{[esp]}? Inside \code{f}, after \code{push ebp; mov ebp, esp}, where is the first argument?
  \item Which registers must a cdecl function preserve, and which may it destroy?
  \item Why is the kernel linked at 1\,MiB rather than at address 0?
  \item What two things does the Multiboot contract explicitly \emph{not} set up for you?
  \item How does a real-mode address like \code{0x07C0:0x0000} map to a physical address?
\end{enumerate}
\end{questionbox}

\section{Going further}

Extend the boot sector: read a key with BIOS \code{int 0x16} and echo it, write a routine that prints a register in
hex, or load a second sector from disk with \code{int 0x13}. That last one is the first step of a custom boot loader.
It is a fun rabbit hole, but only a detour for this project, so set yourself a time limit. For deeper assembly
practice, read the Intel SDM Volume 1 chapter on procedure calls, then write \code{memcpy} and \code{strlen} in NASM
and test them against your C versions from Phase 0.

\nextphase{2}{Getting to C: Screen, Serial and kprintf}
\booklegend
\end{document}
